\documentclass[12pt]{amsart}
\usepackage[margin=1in]{geometry}
\pagestyle{plain}

\usepackage{amsmath,amssymb}
\usepackage{enumerate}
\usepackage[shortlabels]{enumitem}
\setlist[enumerate]{itemsep=4pt, topsep=6pt, parsep=0pt}

\usepackage{graphicx}
\usepackage[T1]{fontenc}
\usepackage[parfill]{parskip}
\renewcommand{\baselinestretch}{1.1}

% Color for correct answers
\usepackage{xcolor}
\definecolor{correctgreen}{rgb}{0,0.6,0}
\definecolor{partialblue}{rgb}{0,0,0.8}
\definecolor{wrongred}{rgb}{0.8,0,0}

\newcommand{\correct}[1]{{\color{correctgreen}\textbf{#1}}}
\newcommand{\partcred}[1]{{\color{partialblue}#1}}
\newcommand{\wrong}[1]{{\color{wrongred}#1}}
\newcommand{\norm}[1]{\left\|#1\right\|}

\title{ESE 2030 QUIZZAM 4 : SOLUTIONS GUIDE\\Fall 2025}
\author{Instructor Solutions}

\begin{document}
\maketitle

\section*{Answer Key Summary}
\begin{center}
\begin{tabular}{|c|c|c|}
\hline
Problem & Correct Answer & Partial Credit \\
\hline
1 & B & -- \\
2 & B & A (small) \\
3 & A & -- \\
4 & A & C (small) \\
5 & C & -- \\
6 & E & D \\
7 & A & B (small) \\
8 & D & -- \\
9 & B & D (large) \\
10 & D & A \\
11 & B & C (almost full) \\
12 & D & C \\
13 & C & -- \\
14 & E & D \\
15 & A & E (small) \\
16 & C & -- \\
17 & D & -- \\
18 & A & -- \\
19 & E & C \\
20 & B & C (small) \\
21 & E & B, C \\
22 & B & -- \\
\hline
\end{tabular}
\end{center}

\textbf{Notes:} This guide provides detailed explanations for each problem. There is often more than one way to understand a given concept, and this guide does not exhaustively cover all approaches. If you have questions, feel free to ask me, a TA, or consult AI tools for clarification.

\newpage

\section*{Detailed Solutions}

\subsection*{Problem 1}
\correct{Answer: (B) $\lambda_2 = -4$}

\textbf{Explanation:} The dominant eigenvalue is the one with largest magnitude (absolute value). Computing magnitudes:
\begin{align*}
|\lambda_1| &= |3| = 3\\
|\lambda_2| &= |-4| = 4\\
|\lambda_3| &= |0| = 0\\
|\lambda_4| = |\lambda_5| &= |1 \pm 2i| = \sqrt{1^2 + 2^2} = \sqrt{5} \approx 2.24
\end{align*}

Since $|\lambda_2| = 4$ is largest, $\lambda_2 = -4$ is the dominant eigenvalue. The sign is irrelevant for dominance—only magnitude matters. The dominant eigenvalue determines the asymptotic rate at which powers $B^k$ grow or decay: here $\|B^k\| \approx 4^k$ for large $k$.

\textbf{Wrong Answers:}
\begin{itemize}
\item (A): $|\lambda_1| = 3 < 4$, so not dominant
\item (C): Zero eigenvalue has smallest magnitude
\item (D): Complex eigenvalues have magnitude $\sqrt{5} \approx 2.24 < 4$
\item (E): There is a clearly dominant eigenvalue
\end{itemize}

\subsection*{Problem 2}
\correct{Answer: (B) $\{e^{2t}, e^{-t}\cos(3t), e^{-t}\sin(3t)\}$}

\partcred{Partial Credit: (A) Mathematically correct but not natural for real ODEs}

\textbf{Explanation:} Complex conjugate eigenvalues $\alpha \pm i\beta$ produce real basis solutions $e^{\alpha t}\cos(\beta t)$ and $e^{\alpha t}\sin(\beta t)$. Here $\alpha = -1$ and $\beta = 3$. The simple real eigenvalue $\lambda = 2$ gives $e^{2t}$. Thus the natural real basis is $\{e^{2t}, e^{-t}\cos(3t), e^{-t}\sin(3t)\}$.

While choice (A) is mathematically valid (complex exponentials are solutions), for real differential equations we typically prefer real-valued basis solutions, which is what (B) provides through Euler's formula.

\textbf{Partial Credit Rationale:} (A) uses complex exponentials which form a valid basis, but not the most natural choice for real ODEs.

\textbf{Wrong Answers:}
\begin{itemize}
\item (C): These would come from a repeated real eigenvalue $\lambda = 2$ with geometric multiplicity 1
\item (D): Incorrect; separates exponential and trigonometric parts
\item (E): Cannot have 4 basis functions for a 3-dimensional solution space
\end{itemize}

\subsection*{Problem 3}
\correct{Answer: (A) $\lambda_1^2 c_1^2 + \cdots + \lambda_n^2 c_n^2$}

\textbf{Explanation:} Since the eigenvectors form an orthonormal basis, we have:
$$A\mathbf{x} = A(c_1\mathbf{q}_1 + \cdots + c_n\mathbf{q}_n) = c_1\lambda_1\mathbf{q}_1 + \cdots + c_n\lambda_n\mathbf{q}_n$$

Because the eigenvectors are orthonormal (a key consequence of the Spectral Theorem), we can compute:
$$\|A\mathbf{x}\|^2 = (c_1\lambda_1)^2 + (c_2\lambda_2)^2 + \cdots + (c_n\lambda_n)^2 = \lambda_1^2 c_1^2 + \cdots + \lambda_n^2 c_n^2$$

This shows how $A$ stretches each eigenspace component independently by the corresponding eigenvalue. The orthonormality is crucial—without it, cross terms would appear.

\textbf{Wrong Answers:}
\begin{itemize}
\item (B): Incorrectly squares the entire sum; would only work if all vectors were parallel
\item (C): Only uses maximum eigenvalue; misses that different components scale differently
\item (D): Forgets to apply $A$ at all; this is just $\|\mathbf{x}\|^2$
\item (E): Uses wrong norm and treats components linearly rather than quadratically
\end{itemize}

\subsection*{Problem 4}
\correct{Answer: (A) I and III only}

\partcred{Partial Credit: (C) Includes II which has wrong factorial}

\textbf{Explanation:} For a Jordan block $J = \lambda I + N$ where $N$ is nilpotent:
$$e^{Jt} = e^{(\lambda I + N)t} = e^{\lambda t}e^{Nt}$$

The key formula is:
$$e^{Nt} = I + Nt + \frac{(Nt)^2}{2!} + \frac{(Nt)^3}{3!} + \cdots$$

For a $2 \times 2$ Jordan block: $e^{Nt} = \begin{bmatrix} 1 & t \\ 0 & 1 \end{bmatrix}$

For a $3 \times 3$ Jordan block: $e^{Nt} = \begin{bmatrix} 1 & t & \frac{t^2}{2} \\ 0 & 1 & t \\ 0 & 0 & 1 \end{bmatrix}$

Checking each matrix:
\begin{itemize}
\item \textbf{I}: Two $2 \times 2$ blocks with $e^{2t}$ on diagonal and $te^{2t}$ on superdiagonal — \textbf{CORRECT} ✓
\item \textbf{II}: One $3 \times 3$ block but has $t^2e^{2t}$ instead of $\frac{t^2}{2}e^{2t}$ — \textbf{WRONG}, missing the factorial
\item \textbf{III}: One $3 \times 3$ block with correct $\frac{t^2}{2}e^{2t}$ term — \textbf{CORRECT} ✓
\item \textbf{IV}: Has $\frac{t^2}{2}e^{2t}$ in wrong position (2,3) instead of (1,3); incorrect block structure
\end{itemize}

\textbf{Partial Credit Rationale:} (C) recognizes that I and III are valid but incorrectly includes II.

\subsection*{Problem 5}
\correct{Answer: (C) The characteristic polynomial of $C$ is $\lambda^3 - 4\lambda^2 + 5\lambda - 2 = 0$}

\textbf{Explanation:} The characteristic polynomial of the companion matrix $C$ is identical to the characteristic polynomial of the differential operator. By substituting $\lambda$ for $D$ in the differential operator $D^3 - 4D^2 + 5D - 2$, we obtain:
$$\lambda^3 - 4\lambda^2 + 5\lambda - 2 = 0$$

The roots of this polynomial are the eigenvalues of $C$, and basis solutions to both formulations (ODE and matrix system) come from these eigenvalues. The conversion to first-order form preserves all solution structure—this is the fundamental connection between higher-order ODEs and matrix systems.

\textbf{Wrong Answers:}
\begin{itemize}
\item (A): Coefficients are not eigenvalues; eigenvalues are roots of the characteristic polynomial
\item (B): Trace equals sum of eigenvalues (which equals 4, the coefficient of $\lambda^2$), not $-1$
\item (D): Companion matrices are not symmetric
\item (E): Both have dimension 3; conversion doesn't change solution space dimension
\end{itemize}

\subsection*{Problem 6}
\correct{Answer: (E) All of them}

\partcred{Partial Credit: (D) Only $\mathcal{M}_1$ and $\mathcal{M}_3$}

\textbf{Explanation:} Real Jordan canonical form consists of:
\begin{itemize}
\item $1 \times 1$ blocks for real eigenvalues: $[\lambda]$
\item $2 \times 2$ blocks for complex conjugate pairs $\alpha \pm i\beta$: $\begin{bmatrix} \alpha & -\beta \\ \beta & \alpha \end{bmatrix}$
\item Jordan blocks with 1's on the superdiagonal for repeated eigenvalues with geometric multiplicity less than algebraic multiplicity
\end{itemize}

Checking each matrix:

\textbf{$\mathcal{M}_1$}: 
\begin{itemize}
\item $[2]$ — real eigenvalue 2
\item $\begin{bmatrix} 1 & -3 \\ 3 & 1 \end{bmatrix}$ — complex pair $1 \pm 3i$
\item $\begin{bmatrix} 4 & 1 & 0 \\ 0 & 4 & 0 \\ 0 & 0 & 4 \end{bmatrix}$ — this is a $2 \times 2$ Jordan block $\begin{bmatrix} 4 & 1 \\ 0 & 4 \end{bmatrix}$ plus a $1 \times 1$ block $[4]$
\end{itemize}
Valid Jordan form! ✓

\textbf{$\mathcal{M}_2$}:
\begin{itemize}
\item $[2]$ — real eigenvalue 2
\item $\begin{bmatrix} 0 & 3 \\ -3 & 0 \end{bmatrix}$ — complex pair $\pm 3i$ (with $\alpha = 0$)
\item $\begin{bmatrix} 4 & 1 \\ 0 & 4 \end{bmatrix}$ — Jordan block for repeated eigenvalue 4
\item $[5]$ — real eigenvalue 5
\end{itemize}
Valid Jordan form! ✓

\textbf{$\mathcal{M}_3$}:
\begin{itemize}
\item $\begin{bmatrix} 1 & -3 \\ 3 & 1 \end{bmatrix}$ — complex pair $1 \pm 3i$
\item $\begin{bmatrix} 4 & 1 & 0 \\ 0 & 4 & 1 \\ 0 & 0 & 4 \end{bmatrix}$ — $3 \times 3$ Jordan block for eigenvalue 4
\item $[2]$ — real eigenvalue 2
\end{itemize}
Valid Jordan form! ✓

\textbf{Partial Credit Rationale:} (D) correctly identifies two matrices but misses that $\mathcal{M}_2$ is also valid.

\subsection*{Problem 7}
\correct{Answer: (A) As $k \to \infty$, the population distribution converges to $\mathbf{v}_1$}

\partcred{Partial Credit: (B) Small — recognizes importance of $\lambda_1 = 1$ but misinterprets its meaning}

\textbf{Explanation:} A stochastic matrix preserves total population (or total probability), which corresponds to the eigenvalue $\lambda_1 = 1$ being dominant. Since $|\lambda_2| = 0.6$ and $|\lambda_3| = 0.2$ are both less than 1, as $k \to \infty$:
$$\mathbf{x}(k) = P^k\mathbf{x}(0) = c_1 \lambda_1^k \mathbf{v}_1 + c_2 \lambda_2^k \mathbf{v}_2 + c_3 \lambda_3^k \mathbf{v}_3 \to c_1 \mathbf{v}_1$$

The eigenvector $\mathbf{v}_1 = \frac{1}{10}\begin{pmatrix} 2 \\ 5 \\ 3 \end{pmatrix}$ indicates the long-term distribution ratios: 20% juvenile, 50% adult, 30% senior. The coefficient $c_1$ is determined by the initial distribution but doesn't affect the ratios.

\textbf{Partial Credit Rationale:} (B) correctly identifies that $\lambda_1 = 1$ is important but incorrectly interprets it to mean the total population stays constant (which is true) rather than recognizing convergence to the stationary distribution.

\textbf{Wrong Answers:}
\begin{itemize}
\item (C): Negative eigenvalues don't make population components negative; they cause oscillating deviations from equilibrium that decay away
\item (D): $|\lambda_2|$ controls convergence rate, not which age group dominates
\item (E): Ignores the dominant eigenvalue $\lambda_1 = 1$; the other eigenvalues affect transients, not the long-term limit
\end{itemize}

\subsection*{Problem 8}
\correct{Answer: (D) The network consists of exactly two separate connected components}

\textbf{Explanation:} A fundamental theorem in spectral graph theory states that the \textbf{multiplicity of the eigenvalue 0} for a graph Laplacian equals the \textbf{number of connected components}. Since $\lambda = 0$ appears with multiplicity 2, the network consists of exactly two separate groups of servers with no connections between the groups.

This is not an error—disconnected networks naturally have multiple zero eigenvalues. The eigenvector for each zero eigenvalue is constant on one component and zero elsewhere, so the eigenspace of 0 has dimension equal to the number of components.

\textbf{Wrong Answers:}
\begin{itemize}
\item (A): Doesn't recognize that multiple zero eigenvalues indicate disconnection
\item (B): Confuses the meaning of $\lambda = 2$; nonzero eigenvalues don't directly count isolated nodes
\item (C): Incorrectly thinks multiple zero eigenvalues indicate an error; this is actually the correct behavior for disconnected graphs
\item (E): Confuses the number of eigenvalues with the number of components; it's the multiplicity of 0 that matters
\end{itemize}

\subsection*{Problem 9}
\correct{Answer: (B) The skew-symmetric part $\beta J$ satisfies $J^2 = -I$, causing the Taylor series to separate into even/odd terms that match the series for $\cos(\beta t)$ and $\sin(\beta t)$}

\partcred{Partial Credit: (D) Large — recognizes connection to Euler's formula}

\textbf{Explanation:} Let $N = \beta J = \begin{bmatrix} 0 & -\beta \\ \beta & 0 \end{bmatrix}$. The key property is:
$$N^2 = \begin{bmatrix} 0 & -\beta \\ \beta & 0 \end{bmatrix}^2 = \begin{bmatrix} -\beta^2 & 0 \\ 0 & -\beta^2 \end{bmatrix} = -\beta^2 I$$

This gives us: $N^3 = -\beta^2 N$, $N^4 = \beta^4 I$, etc.

When we expand $e^{Nt}$ as a Taylor series:
$$e^{Nt} = I + Nt + \frac{(Nt)^2}{2!} + \frac{(Nt)^3}{3!} + \frac{(Nt)^4}{4!} + \cdots$$

The even powers give:
$$I\left(1 - \frac{\beta^2t^2}{2!} + \frac{\beta^4t^4}{4!} - \cdots\right) = \cos(\beta t) \cdot I$$

The odd powers give:
$$N\left(t - \frac{\beta^2t^3}{3!} + \frac{\beta^4t^5}{5!} - \cdots\right) = \frac{1}{\beta}N \cdot \beta t\left(1 - \frac{(\beta t)^2}{3!} + \cdots\right) = \frac{1}{\beta}N\sin(\beta t)$$

The appearance of sines and cosines comes from how their Taylor series naturally emerge when computing matrix powers. The crucial property is $N^2 = -\beta^2 I$, which causes the alternating signs in the series expansion.

\textbf{Partial Credit Rationale:} (D) recognizes the connection to Euler's formula, which is related but doesn't explain the mechanism from first principles.

\textbf{Wrong Answers:}
\begin{itemize}
\item (A): Describes the phenomenon but doesn't explain the mechanism; this is "what" not "why"
\item (C): Confuses cause and effect—$J$ represents rotation BECAUSE the exponential has this form, not vice versa
\item (E): The determinant property is irrelevant to why trigonometric functions appear in the series expansion
\end{itemize}

\subsection*{Problem 10}
\correct{Answer: (D) The sequence converges to the direction of $\mathbf{v}_2$, but oscillating signs}

\partcred{Partial Credit: (A) Recognizes convergence to zero but misses dominant direction and oscillation}

\textbf{Explanation:} Since the eigenvectors are linearly independent:
$$\mathbf{x}_k = A^k\mathbf{x}_0 = 2(0.6)^k\mathbf{v}_1 + 3(-0.8)^k\mathbf{v}_2 + 5(0.4)^k\mathbf{v}_3$$

The dominant eigenvalue (largest magnitude) is $\lambda_2 = -0.8$ with $|\lambda_2| = 0.8$. The other eigenvalues have smaller magnitudes: $|\lambda_1| = 0.6$ and $|\lambda_3| = 0.4$.

As $k \to \infty$:
\begin{itemize}
\item The ratios $|\lambda_1|/|\lambda_2| = 0.75$ and $|\lambda_3|/|\lambda_2| = 0.5$ mean the other components decay much faster relative to the $\mathbf{v}_2$ component
\item Therefore $\mathbf{x}_k \approx 3(-0.8)^k\mathbf{v}_2$ for large $k$
\item The normalized sequence $\mathbf{x}_k/\|\mathbf{x}_k\|$ oscillates between $+\mathbf{v}_2/\|\mathbf{v}_2\|$ and $-\mathbf{v}_2/\|\mathbf{v}_2\|$ (alternating signs due to the negative eigenvalue)
\item The magnitude $\|\mathbf{x}_k\|$ decays to zero since $|\lambda_2| = 0.8 < 1$
\end{itemize}

The sequence converges to zero along the direction of $\mathbf{v}_2$, with oscillating sign.

\textbf{Partial Credit Rationale:} (A) correctly recognizes that all eigenvalue magnitudes are less than 1, so the sequence converges to zero, but misses the key insight about the dominant eigenvector direction and oscillatory behavior.

\textbf{Wrong Answers:}
\begin{itemize}
\item (B): Incorrectly identifies $\lambda_1 = 0.6$ as dominant; dominance is determined by magnitude
\item (C): All eigenvalues have magnitude less than 1, so the sequence decays
\item (E): While eigencomponents act independently, this doesn't describe the asymptotic behavior where one component dominates
\end{itemize}

\subsection*{Problem 11}
\correct{Answer: (B) $-A$ has a dominant real negative eigenvalue with positive eigenvector}

\partcred{Partial Credit: (C) Almost full — gets eigenvalue sign correct but not eigenvector sign}

\textbf{Explanation:} If $\mathbf{v}$ is an eigenvector of $A$ with eigenvalue $\lambda$, then $A\mathbf{v} = \lambda\mathbf{v}$, which means:
$$(-A)\mathbf{v} = -\lambda\mathbf{v}$$

Thus $\mathbf{v}$ is also an eigenvector of $-A$ with eigenvalue $-\lambda$. The key insight: \textbf{negating a matrix negates its eigenvalues but preserves eigenvectors}.

Since $A$ has dominant eigenvalue $\lambda_* = \rho(A) > 0$ with positive eigenvector $\mathbf{v}_*$ (by Perron-Frobenius), the matrix $-A$ has:
\begin{itemize}
\item Eigenvalue $-\lambda_* < 0$ (negative)
\item With the SAME positive eigenvector $\mathbf{v}_*$
\item This eigenvalue has magnitude $|-\lambda_*| = \lambda_* = \rho(A)$, making it dominant
\end{itemize}

\textbf{Partial Credit Rationale:} (C) correctly identifies the negative eigenvalue but incorrectly assumes the eigenvector changes sign. Eigenvectors are unchanged by negating the matrix.

\textbf{Wrong Answers:}
\begin{itemize}
\item (A): Negating the matrix negates the eigenvalues
\item (D): While Perron-Frobenius doesn't directly apply to matrices with all negative entries, we can determine the eigenstructure by understanding how negation affects eigenvalues and eigenvectors
\item (E): (B) is entirely accurate
\end{itemize}

\subsection*{Problem 12}
\correct{Answer: (D) $R_k Q_k$ is similar to $A_k$ via $A_{k+1} = Q_k^T A_k Q_k$, preserving eigenvalues across iterations}

\partcred{Partial Credit: (C) Recognizes that $Q_k R_k = A_k$ but misses the crucial similarity relationship}

\textbf{Explanation:} The key insight is that:
$$R_k Q_k = Q_k^T Q_k R_k Q_k = Q_k^T (Q_k R_k) Q_k = Q_k^T A_k Q_k$$

since $A_k = Q_k R_k$ and $Q_k^T Q_k = I$ (orthogonality). This shows $A_{k+1}$ is similar to $A_k$ via an orthogonal similarity transformation, guaranteeing they have identical eigenvalues.

The sequence of similarity transformations:
$$A_0 \to A_1 \to A_2 \to \cdots$$
gradually reveals the eigenvalue structure (converging to triangular or block-triangular form) while preserving the spectrum at each step.

Choice (C) is tempting but imprecise—while $Q_k R_k = A_k$ is true, the deeper reason we reverse the order is to create the similarity transformation that preserves eigenvalues while changing the matrix form toward something more revealing.

\textbf{Partial Credit Rationale:} (C) recognizes that $Q_k R_k = A_k$ but misses the crucial similarity relationship that makes the algorithm converge to useful form.

\textbf{Wrong Answers:}
\begin{itemize}
\item (A): $R_k Q_k$ is NOT generally upper triangular; the algorithm converges to triangular form over many iterations, not in one step
\item (B): While matrix multiplication is non-commutative, this observation doesn't explain why the algorithm works
\item (E): $R_k Q_k$ does not generally preserve symmetry unless additional structure is present
\end{itemize}

\subsection*{Problem 13}
\correct{Answer: (C) Matrix $B$ converges faster}

\textbf{Explanation:} The convergence rate of the power method is determined by the ratio $|\lambda_2|/|\lambda_1|$ where $\lambda_1$ is the dominant eigenvalue and $\lambda_2$ is the second-largest in magnitude.

For matrix $A$: ratio is $|9|/|10| = 0.9$

For matrix $B$: ratio is $|2|/|10| = 0.2$

The smaller this ratio, the faster the convergence. Matrix $B$ has ratio $0.2 < 0.9$, so it converges much faster.

\textbf{Intuitively}: When the dominant eigenvalue is well-separated from the others (as in $B$), the $\mathbf{v}_1$ component grows much faster relative to other components, leading to rapid convergence. When eigenvalues are close (as in $A$), the $\mathbf{v}_2$ component remains significant for many iterations, slowing convergence.

After $k$ iterations, the relative error is approximately $(|\lambda_2|/|\lambda_1|)^k$. For $A$: $(0.9)^k$ vs. for $B$: $(0.2)^k$ — the latter decays exponentially faster.

\textbf{Wrong Answers:}
\begin{itemize}
\item (A): Doesn't recognize that separation matters
\item (B): Backward reasoning; larger second eigenvalue means slower convergence
\item (D): While initial condition affects early iterations, the asymptotic convergence rate is determined by the eigenvalue ratio
\item (E): The power method doesn't require Perron-Frobenius; it applies to any diagonalizable matrix
\end{itemize}

\subsection*{Problem 14}
\correct{Answer: (E) Need to know the dominant eigenvector to determine this}

\partcred{Partial Credit: (D) Recognizes need to compute limit but doesn't identify the eigenvector method}

\textbf{Explanation:} Stochastic matrices (where columns sum to 1 and entries are non-negative) always have $\lambda = 1$ as an eigenvalue. The corresponding eigenvector, normalized so its entries sum to 1, gives the \textbf{stationary distribution}—the long-term market share that the system converges to regardless of initial distribution.

After many iterations:
$$\mathbf{x}_k = P^k\mathbf{x}_0 \to \text{stationary distribution} = \text{eigenvector for } \lambda = 1$$

This is a fundamental result in Markov chain theory: the eigenvalue $\lambda = 1$ is dominant (by Perron-Frobenius for irreducible stochastic matrices), and its positive eigenvector describes the equilibrium.

To find the actual market shares, we would need to:
\begin{enumerate}
\item Solve $(P - I)\mathbf{v} = \mathbf{0}$ for the eigenvector
\item Normalize so entries sum to 1
\end{enumerate}

Without doing this calculation, we cannot determine the specific percentages.

\textbf{Partial Credit Rationale:} (D) correctly recognizes that we need to compute a limit, which is related to finding the eigenvector.

\textbf{Wrong Answers:}
\begin{itemize}
\item (A): Initial distribution doesn't determine long-term behavior
\item (B), (C): Random guesses without calculation
\end{itemize}

\subsection*{Problem 15}
\correct{Answer: (A) Over time, the robot spends the same time in room 1 as in room 4}

\partcred{Partial Credit: (E) True but not via symmetry — may be true for other reasons}

\textbf{Explanation:} The stationary distribution of a random walk on a graph reflects the long-term proportion of time spent in each state. When two rooms have \textbf{identical rows AND columns} in the transition matrix, they are \textbf{topologically equivalent}—they have the same connectivity structure to the rest of the network.

This symmetry forces the stationary distribution to assign equal probability to these rooms. Since the robot's long-term behavior is described by the eigenvector for $\lambda = 1$, and this eigenvector must respect the symmetry of $P$, rooms 1 and 4 must have equal components.

This is a manifestation of the symmetry principle in Markov chains: states with identical transition patterns have identical stationary probabilities.

Verification: Looking at the matrix, rooms 1 and 4 both:
\begin{itemize}
\item Have 2 neighbors (degree 2)
\item Connect to the same set of rooms (2 and 3)
\item Have identical transition probabilities
\end{itemize}

\textbf{Partial Credit Rationale:} (E) may be true (rooms 2 and 3 have degree 3 vs. rooms 1 and 4 have degree 2, so may have higher stationary probability), but this doesn't follow from the symmetry argument.

\textbf{Wrong Answers:}
\begin{itemize}
\item (B): No such symmetry exists between rooms 2 and 3
\item (C): Instantaneous probabilities don't determine long-term behavior
\item (D): Eigenvalues are properties of the matrix, not individual states
\end{itemize}

\subsection*{Problem 16}
\correct{Answer: (C) $C$ has eigenvalue $\lambda = 3$ appearing twice in the characteristic polynomial, with a 1-dimensional eigenspace}

\textbf{Explanation:} The general solution structure reveals the eigenvalue structure:
\begin{itemize}
\item $e^{-2t}$ indicates eigenvalue $\lambda = -2$ appearing once
\item $e^{3t}$ and $te^{3t}$ indicate eigenvalue $\lambda = 3$ appearing twice
\end{itemize}

The presence of BOTH $e^{3t}$ and $te^{3t}$ is the key diagnostic. The polynomial term $t$ always signals a deficiency in the number of eigenvectors relative to how many times the eigenvalue appears (geometric multiplicity < algebraic multiplicity).

If the eigenspace for $\lambda = 3$ were 2-dimensional, we would have two independent eigenvectors for $\lambda = 3$, producing two independent exponential solutions of the form $e^{3t}$ (in different directions). Instead, the $te^{3t}$ term arises precisely because there is only ONE eigenvector for $\lambda = 3$, requiring a generalized eigenvector to complete the basis.

Summary for $\lambda = 3$:
\begin{itemize}
\item Algebraic multiplicity: 2 (appears twice in characteristic polynomial)
\item Geometric multiplicity: 1 (one eigenvector, one generalized eigenvector)
\item Basis solutions: $e^{3t}$ (from eigenvector), $te^{3t}$ (from generalized eigenvector)
\end{itemize}

\textbf{Wrong Answers:}
\begin{itemize}
\item (A): Misses that the polynomial term signals eigenvalue repetition
\item (B): Backwards; a 2-dimensional eigenspace would give two independent exponential solutions, not the polynomial term
\item (D): Contradictory reasoning; a 2-dimensional eigenspace means two independent eigenvectors
\item (E): Complex eigenvalues produce trigonometric functions like $\cos(\beta t)$ and $\sin(\beta t)$, not polynomial-exponential products
\end{itemize}

\subsection*{Problem 17}
\correct{Answer: (D) Block upper triangular: three $1 \times 1$ blocks and one $2 \times 2$ block}

\textbf{Explanation:} The QR algorithm converges to \textbf{real Schur form}: block upper triangular with:
\begin{itemize}
\item $1 \times 1$ blocks for real eigenvalues (regardless of multiplicity)
\item $2 \times 2$ blocks for complex conjugate pairs
\end{itemize}

For this matrix:
\begin{itemize}
\item Three real eigenvalues: $3, 3, -1$ → three $1 \times 1$ blocks
\item One complex conjugate pair: $1 \pm 2i$ → one $2 \times 2$ block
\end{itemize}

The $2 \times 2$ block has the form $\begin{bmatrix} 1 & -2 \\ 2 & 1 \end{bmatrix}$ (or similar, depending on eigenvector structure).

Important: The QR algorithm:
\begin{itemize}
\item Does NOT produce Jordan form (would need generalized eigenvectors)
\item Does NOT produce complex entries (real arithmetic throughout)
\item Is NOT prevented from converging by repeated eigenvalues
\item Does NOT produce a diagonal matrix unless all eigenvalues are real and distinct
\end{itemize}

\textbf{Wrong Answers:}
\begin{itemize}
\item (A): Real sequence cannot have complex entries
\item (B): QR gives Schur form, not Jordan form
\item (C): Complex eigenvalues require $2 \times 2$ blocks in real arithmetic
\item (E): Repeated eigenvalues do not prevent convergence
\end{itemize}

\subsection*{Problem 18}
\correct{Answer: (A) There are three positive eigenvalues}

\textbf{Explanation:} The matrix $A = X^TX$ is a \textbf{Gram matrix}. When $X$ has linearly independent columns:
\begin{itemize}
\item $A$ is symmetric: $(X^TX)^T = X^T(X^T)^T = X^TX$
\item $A$ is positive definite: $\mathbf{v}^T A \mathbf{v} = \mathbf{v}^T X^T X \mathbf{v} = \|X\mathbf{v}\|^2 > 0$ for all $\mathbf{v} \neq \mathbf{0}$
\end{itemize}

By the Spectral Theorem, $A$ can be orthogonally diagonalized with real eigenvalues. The positive definiteness ensures all eigenvalues are positive.

Since $X$ is $5 \times 3$, the matrix $A = X^TX$ is $3 \times 3$. Therefore, $A$ has exactly three eigenvalues, all positive.

\textbf{Wrong Answers:}
\begin{itemize}
\item (B): $A$ is $3 \times 3$ and full rank (positive definite), so no zero eigenvalues
\item (C): Positive definiteness doesn't prevent repeated eigenvalues
\item (D): Wrong dimension; $A$ is $3 \times 3$, not $5 \times 5$
\item (E): Linear independence of columns guarantees positive definiteness
\end{itemize}

\subsection*{Problem 19}
\correct{Answer: (E) All three $\mathbf{v}_i$ exist and are mutually orthogonal}

\partcred{Partial Credit: (C) Recognizes orthogonality can be achieved but doesn't appreciate that Spectral Theorem guarantees this}

\textbf{Explanation:} The \textbf{Spectral Theorem} guarantees that every symmetric matrix has an orthonormal basis of eigenvectors, regardless of whether eigenvalues are repeated.

For this matrix:
\begin{enumerate}
\item Symmetric matrices are always diagonalizable, so geometric multiplicity equals algebraic multiplicity for each eigenvalue. The eigenspace for $\lambda = 3$ is 2-dimensional.
\item The theorem guarantees not just that eigenvectors from different eigenspaces are orthogonal, but that we can choose an ENTIRE orthonormal eigenbasis.
\item This means $\mathbf{v}_1$, $\mathbf{v}_2$, and $\mathbf{v}_3$ can all be chosen mutually orthogonal.
\end{enumerate}

While it's true we might use Gram-Schmidt within the eigenspace for $\lambda = 3$ to construct such a basis, the Spectral Theorem guarantees this construction will succeed—this is the key difference from general (non-symmetric) matrices where repeated eigenvalues might have deficient eigenspaces requiring generalized eigenvectors.

\textbf{Partial Credit Rationale:} (C) recognizes that all three eigenvectors exist and can be made orthogonal via Gram-Schmidt, but doesn't fully appreciate that the Spectral Theorem guarantees this is always possible for symmetric matrices.

\textbf{Wrong Answers:}
\begin{itemize}
\item (A): Confuses symmetric matrices with general matrices; symmetric matrices always have full-dimensional eigenspaces
\item (B): True for eigenvectors from different eigenspaces, but misses that orthogonality is guaranteed within eigenspaces too
\item (D): Repeated eigenvalues can cause non-diagonalizability for general matrices, but never for symmetric matrices
\end{itemize}

\subsection*{Problem 20}
\correct{Answer: (B) $\rho(A) = 3$}

\partcred{Partial Credit: (C) Small — correctly computes magnitude of complex eigenvalues but selects wrong maximum}

\textbf{Explanation:} The spectral radius is the maximum magnitude among all eigenvalues:
$$\rho(A) = \max\{|\lambda_i|\}$$

Computing magnitudes:
\begin{align*}
|\lambda_1| &= |2| = 2\\
|\lambda_2| &= |-3| = 3\\
|\lambda_3| &= |1+2i| = \sqrt{1^2 + 2^2} = \sqrt{5} \approx 2.24\\
|\lambda_4| &= |1-2i| = \sqrt{1^2 + (-2)^2} = \sqrt{5} \approx 2.24
\end{align*}

The maximum is $|\lambda_2| = 3$, so $\rho(A) = 3$.

Note: The spectral radius is always a non-negative real number, even when the matrix has complex eigenvalues. It determines the asymptotic growth rate of matrix powers: $\|A^k\| \approx \rho(A)^k$ for large $k$.

\textbf{Partial Credit Rationale:} (C) correctly computes $|\lambda_3| = |\lambda_4| = \sqrt{5}$ but doesn't compare all magnitudes to find the maximum.

\textbf{Wrong Answers:}
\begin{itemize}
\item (A): This is $\sqrt{3}$, which doesn't correspond to any eigenvalue magnitude
\item (D): This would be $\rho(A)^2$ if $\rho(A) = \sqrt{5}$, but spectral radius is not squared
\item (E): The spectral radius is clearly 3
\end{itemize}

\subsection*{Problem 21}
\correct{Answer: (E) Both (B) and (C)}

\partcred{Partial Credit: (B) or (C) individually}

\textbf{Explanation:} For diagonalizability, geometric multiplicity must equal algebraic multiplicity for each eigenvalue.

For $\lambda = 2$:
\begin{itemize}
\item Algebraic multiplicity: 3 (from $(\ lambda - 2)^3$)
\item Geometric multiplicity: $\dim(\ker(A - 2I)) = 2$
\item Mismatch: $2 < 3$, so NOT diagonalizable
\end{itemize}

For $\lambda = -1$:
\begin{itemize}
\item Algebraic multiplicity: 2 (from $(\lambda + 1)^2$)
\item Geometric multiplicity: $\dim(\ker(A + I)) = 2$
\item Match: $2 = 2$, so this eigenvalue is fine
\end{itemize}

Since there is a deficiency for $\lambda = 2$, the matrix $A$ is NOT diagonalizable.

Furthermore, the deficiency $(3-2=1)$ tells us there must be at least one Jordan block of size at least 2 for eigenvalue 2. Specifically, the Jordan form will have either:
\begin{itemize}
\item One $2 \times 2$ block and one $1 \times 1$ block for $\lambda = 2$, OR
\item One $3 \times 3$ block for $\lambda = 2$
\end{itemize}

(The exact structure depends on $\dim(\ker(A-2I)^2)$, but we know at least one block has size $\geq 2$.)

Additional verification for (D): By rank-nullity, $\text{rank}(A - 2I) = 5 - \dim(\ker(A-2I)) = 5 - 2 = 3$, not 2.

\textbf{Partial Credit Rationale:} (B) or (C) individually correctly identify one consequence of the multiplicity mismatch.

\textbf{Wrong Answers:}
\begin{itemize}
\item (A): Directly contradicts the multiplicity mismatch
\item (D): Rank-nullity gives rank = 3, not 2
\end{itemize}

\subsection*{Problem 22}
\correct{Answer: (B) Two complex conjugate eigenvalues with negative real part}

\textbf{Explanation:} Damped oscillations require two features in the solution behavior:
\begin{enumerate}
\item \textbf{Oscillatory motion} — comes from complex conjugate eigenvalues
\item \textbf{Decay over time} — comes from negative real parts
\end{enumerate}

Complex conjugate eigenvalues $\lambda = \alpha \pm i\beta$ (with $\beta \neq 0$) produce solutions of the form:
$$e^{\alpha t}(A\cos(\beta t) + B\sin(\beta t))$$

For damped oscillations:
\begin{itemize}
\item The trigonometric functions $\cos(\beta t)$ and $\sin(\beta t)$ provide the oscillation
\item The exponential factor $e^{\alpha t}$ controls growth/decay
\item When $\alpha < 0$ (negative real part), we get $e^{\alpha t} \to 0$ as $t \to \infty$, causing the oscillations to decay
\end{itemize}

Therefore, damped oscillations correspond to complex eigenvalues with negative real part.

\textbf{Wrong Answers:}
\begin{itemize}
\item (A): Two distinct real negative eigenvalues give exponential decay without oscillation (overdamped)
\item (C): Repeated real negative eigenvalue gives fastest decay without oscillation (critically damped)
\item (D): Positive real eigenvalues cause unbounded exponential growth, not damping
\item (E): Complex eigenvalues with positive real part produce growing oscillations, not damped ones
\end{itemize}

\end{document}